\documentclass[12pt]{article}
\usepackage[utf8]{inputenc}
\usepackage[T1]{fontenc}
\usepackage[margin=1in]{geometry}
\usepackage{setspace}
\usepackage{xcolor}
\usepackage[colorlinks=true, urlcolor=blue!60!black, linkcolor=blue!60!black]{hyperref}
\setlength{\parindent}{2em}
\setlength{\parskip}{1.25em}
\renewcommand{\baselinestretch}{1.0}
\begin{document}
\begin{center}
{\LARGE \textit{Lujan v. Defenders of Wildlife}, 504 U.S. 555 (1992)\par}
\vspace{0.8em}
{\large Case Brief\par}
\end{center}
\vspace{1em}

\section{Memory Jogger}
Establishing the modern three-part Article III standing test, the Supreme Court held that environmental plaintiffs lacked an ``actual or imminent'' injury in fact based on vague ``some day'' plans to revisit foreign wildlife habitats and that Congress cannot confer standing through a citizen-suit provision to vindicate a generalized procedural grievance (\href{https://www.courtlistener.com/opinion/112747/lujan-v-defenders-of-wildlife/}{CourtListener}; \href{https://supreme.justia.com/cases/federal/us/504/555/}{Justia}).

\section{Detailed Case Facts}
Section 7(a)(2) of the Endangered Species Act of 1973 (ESA), 16 U.S.C. \S~1536(a)(2), requires each federal agency, in consultation with the Secretary of the Interior (or Commerce, as appropriate), to ensure that any action ``authorized, funded, or carried out by such agency'' is not likely to jeopardize the continued existence of any endangered or threatened species or destroy or adversely modify critical habitat (\href{https://www.law.cornell.edu/supct/html/90-1424.ZO.html}{Cornell LII}; \href{https://supreme.justia.com/cases/federal/us/504/555/}{504 U.S. at 558}). In 1978, the Fish and Wildlife Service (FWS) and the National Marine Fisheries Service (NMFS) promulgated a joint regulation extending \S~7(a)(2)'s consultation mandate to federal actions taken in foreign nations (43 Fed. Reg. 874). After reexamining that position, the Secretaries proposed a revised rule in 1983 (48 Fed. Reg. 29990) and promulgated a final joint regulation in 1986 (51 Fed. Reg. 19926; 50 C.F.R. \S~402.01) limiting \S~7(a)(2)'s geographic scope to actions within the United States or on the high seas (\href{https://supreme.justia.com/cases/federal/us/504/555/}{504 U.S. at 558--59}).

Defenders of Wildlife and other wildlife conservation organizations sued the Secretary of the Interior (originally Donald Hodel, later Manuel Lujan, Jr.) under the ESA's citizen-suit provision, 16 U.S.C. \S~1540(g), seeking a declaratory judgment that the 1986 regulation violated the ESA and an injunction requiring the Secretary to restore the 1978 extraterritorial interpretation (\href{https://supreme.justia.com/cases/federal/us/504/555/}{504 U.S. at 559, 571--72}). To establish standing at summary judgment, respondents relied primarily on affidavits from two members: Joyce Kelly, who had traveled to Egypt in 1986, observed the habitat of the endangered Nile crocodile, and stated that she intended to return to observe the crocodile and would be harmed by U.S. involvement (via the Bureau of Reclamation) in rehabilitating the Aswan High Dam; and Amy Skilbred, who had visited Sri Lanka in 1981, observed habitat of the endangered Asian elephant and leopard near the site of the U.S. Agency for International Development (AID)-funded Mahaweli project without seeing the animals, and stated she intended to return ``in the future''---though she admitted in deposition that she had no current plans to return because of Sri Lanka's civil war (\href{https://supreme.justia.com/cases/federal/us/504/555/}{504 U.S. at 563--64}).

\section{Procedural History}
The U.S. District Court for the District of Minnesota initially granted the Secretary's motion to dismiss for lack of standing (\textit{Defenders of Wildlife v. Hodel}, 658 F. Supp. 43 (D. Minn. 1987)) (\href{https://supreme.justia.com/cases/federal/us/504/555/}{504 U.S. at 559}). A divided panel of the U.S. Court of Appeals for the Eighth Circuit reversed, holding that respondents had standing (851 F.2d 1035 (8th Cir. 1988)). On remand, on cross-motions for summary judgment, the District Court denied the Secretary's renewed standing motion (treating the Eighth Circuit's prior ruling as controlling), granted summary judgment to respondents on the merits, and ordered the Secretary to publish a revised regulation (707 F. Supp. 1082 (D. Minn. 1989)). The Eighth Circuit affirmed (911 F.2d 117 (8th Cir. 1990)), holding both that respondents established standing (including via a ``procedural injury'' under the ESA's citizen-suit provision) and that \S~7(a)(2) applied abroad (\href{https://supreme.justia.com/cases/federal/us/504/555/}{504 U.S. at 559, 571--72}). The Supreme Court granted certiorari (500 U.S. 915 (1991)), heard oral argument on December 3, 1991, and reversed and remanded on June 12, 1992 (\href{https://www.oyez.org/cases/1991/90-1424}{Oyez}; \href{https://supreme.justia.com/cases/federal/us/504/555/}{504 U.S. at 555}).

\section{Judicial Votes}
\textbf{Vote: 7--2 in the judgment} (reversing the Eighth Circuit), with a \textbf{6--3 majority} holding that respondents lacked Article III standing and a \textbf{4-Justice plurality} on redressability (\href{https://supreme.justia.com/cases/federal/us/504/555/}{504 U.S. at 556--57}; \href{https://tile.loc.gov/storage-services/service/ll/usrep/usrep504/usrep504555/usrep504555.pdf}{U.S. Reports}).
\begin{itemize}
\item \textbf{Opinion of the Court (Parts I, II, III-A, and IV) (6 Justices):} Justice Antonin Scalia delivered the opinion of the Court, joined by Chief Justice William H. Rehnquist and Justices Byron R. White, Anthony M. Kennedy, David H. Souter, and Clarence Thomas.
\item \textbf{Plurality Opinion (Part III-B, Redressability) (4 Justices):} Justice Antonin Scalia, joined by Chief Justice William H. Rehnquist and Justices Byron R. White and Clarence Thomas.
\item \textbf{Concurring in Part and Concurring in the Judgment (2 Justices):} Justice Anthony M. Kennedy, joined by Justice David H. Souter (joined Parts I, II, III-A, and IV, but declined to join Part III-B).
\item \textbf{Concurring in the Judgment (1 Justice):} Justice John Paul Stevens (concluded respondents had standing, but concurred in reversal on the merits because \S~7(a)(2) does not apply abroad).
\item \textbf{Dissenting (2 Justices):} Justice Harry A. Blackmun, joined by Justice Sandra Day O'Connor.
\end{itemize}

\section{Holding}
\textbf{Judgment:} Reversed and remanded; the Court of Appeals erred in denying the United States' motion for summary judgment for lack of standing (\href{https://supreme.justia.com/cases/federal/us/504/555/}{504 U.S. at 578}).

\textbf{Rule Announced:} The ``irreducible constitutional minimum'' of Article III standing requires a plaintiff to establish three elements: (1) an \textit{injury in fact} that is concrete and particularized, and actual or imminent rather than conjectural or hypothetical; (2) a \textit{causal connection} (fairly traceable to the defendant's challenged conduct rather than independent third-party action); and (3) \textit{redressability} (likely, not merely speculative, that a favorable judicial decision will redress the injury) (\href{https://www.law.cornell.edu/constitution-conan/article-3/section-2/clause-1/overview-of-the-lujan-test}{Constitution Annotated}; \href{https://supreme.justia.com/cases/federal/us/504/555/}{504 U.S. at 560--61}). Each element must be supported with the degree of evidence required at each stage of litigation (specific facts by affidavit or evidence at summary judgment). Furthermore, Congress cannot confer Article III standing through a citizen-suit provision authorizing ``any person'' to sue to enforce an abstract procedural duty unconnected to a concrete, particularized injury of the plaintiff (\href{https://supreme.justia.com/cases/federal/us/504/555/}{504 U.S. at 572--78}).

\textbf{Holding vs. Dicta:} The binding majority holdings (Parts II, III-A, and IV) are the three-part Article III standing framework, the rule that indefinite ``some day'' travel intentions and broad ``ecosystem,'' ``animal,'' or ``vocational'' nexus theories fail to establish imminent injury in fact at summary judgment, and the holding that Article III bars Congress from converting a generalized public grievance over agency procedure into a justiciable individual right. By contrast, Part III-B's redressability analysis (concluding that nonparty funding agencies would not be bound by the Secretary's revised regulation and supplied only a small fraction of project funding) commanded only a four-justice plurality and is therefore not binding precedent (\href{https://supreme.justia.com/cases/federal/us/504/555/}{504 U.S. at 556, 568--71}). Likewise, Footnote 7's statement that plaintiffs enforcing procedural rights linked to a concrete interest (such as a neighbor challenging a dam's missing environmental impact statement) need not meet ``all the normal standards for redressability and immediacy'' is influential dicta that has nonetheless shaped subsequent environmental standing doctrine (\href{https://supreme.justia.com/cases/federal/us/504/555/}{504 U.S. at 572 n.7}).

\section{Analysis of Opinions}

\subsection{Majority Opinion (Parts I, II, III-A, and IV) and Plurality Opinion (Part III-B) (Justice Scalia)}
In Part II, Justice Scalia anchored standing squarely in the Constitution's separation of powers, explaining that Article III's ``Cases'' and ``Controversies'' limitation defines the proper judicial role (\href{https://supreme.justia.com/cases/federal/us/504/555/}{504 U.S. at 559--60}). He synthesized the three-element standing test (injury in fact, traceability, redressability), emphasized that plaintiffs must prove standing with ``specific facts'' at summary judgment under Fed. R. Civ. P. 56(e), and stressed that standing is ``substantially more difficult'' to establish when the plaintiff is not the direct object of the challenged government regulation or inaction (\href{https://supreme.justia.com/cases/federal/us/504/555/}{504 U.S. at 560--62}).

In Part III-A, while acknowledging that observing wildlife for aesthetic purposes is a cognizable interest under \textit{Sierra Club v. Morton}, 405 U.S. 727 (1972), the Court held that Kelly's and Skilbred's affidavits failed to show ``imminent'' injury. Past visits proved nothing without continuing adverse effects, and ``some day'' intentions to return---without concrete plans or any specification of when---could not satisfy Article III (\href{https://supreme.justia.com/cases/federal/us/504/555/}{504 U.S. at 563--64}). The Court also rejected respondents' three ``nexus'' theories: ``ecosystem nexus'' conflicted with \textit{Lujan v. National Wildlife Federation}, 497 U.S. 871 (1990), which required use of the specific area affected; and ``animal nexus'' and ``vocational nexus'' stretched standing ``into pure speculation and fantasy'' by allowing a Bronx Zoo visitor or elephant keeper to challenge an AID project in Sri Lanka (\href{https://supreme.justia.com/cases/federal/us/504/555/}{504 U.S. at 565--67}).

In Part III-B (plurality only), Scalia argued respondents also failed redressability for two reasons: (1) the funding agencies (such as AID) were not parties to the suit and contested whether the Secretary of the Interior's consultation regulation bound them; and (2) U.S. agencies supplied only a fraction of the foreign projects' funding (AID provided less than 10\% of the Mahaweli project), leaving it conjectural whether ending U.S. funding would halt or mitigate harm to endangered species (\href{https://supreme.justia.com/cases/federal/us/504/555/}{504 U.S. at 568--71}).

In Part IV, the majority rejected the Eighth Circuit's ``procedural injury'' holding under the ESA's citizen-suit clause, 16 U.S.C. \S~1540(g). Reviewing generalized-grievance precedents from \textit{Fairchild v. Hughes} (1922) and \textit{Massachusetts v. Mellon} (1923) to \textit{United States v. Richardson} (1974) and \textit{Allen v. Wright} (1984), Scalia concluded that Congress cannot override Article III by conferring an abstract, self-contained procedural right on all citizens. Allowing Congress to convert the public interest in lawful administration into an individual right would transfer the President's Article II, \S~3 duty to ``take Care that the Laws be faithfully executed'' from the Executive to the courts (\href{https://supreme.justia.com/cases/federal/us/504/555/}{504 U.S. at 571--78}).

\subsection{Concurrence in Part and in the Judgment (Justice Kennedy, joined by Justice Souter)}
Justice Kennedy joined Parts I, II, III-A, and IV, agreeing that Kelly and Skilbred failed to show concrete, imminent injury and that the ESA's citizen-suit clause could not create standing without a concrete injury (\href{https://supreme.justia.com/cases/federal/us/504/555/}{504 U.S. at 579--81}). Crucially, however, Kennedy declined to join Part III-B on redressability and wrote separately to temper two aspects of Scalia's opinion: first, he was unwilling to foreclose the possibility that a nexus theory might succeed on a different record; second, he emphasized that ``Congress has the power to define injuries and articulate chains of causation that will give rise to a case or controversy where none existed before,'' provided Congress at a minimum identifies the injury it seeks to vindicate and relates that injury to the class of persons entitled to sue (\href{https://supreme.justia.com/cases/federal/us/504/555/}{504 U.S. at 579--80}).

\subsection{Concurrence in the Judgment (Justice Stevens)}
Justice Stevens disagreed with the Court on standing, arguing that the imminence of an aesthetic or scientific injury should be measured by the timing and likelihood of the threatened environmental harm rather than when the plaintiff plans to visit, and that Kelly and Skilbred's past visits and professional backgrounds showed a genuine interest sufficient to survive summary judgment (\href{https://supreme.justia.com/cases/federal/us/504/555/}{504 U.S. at 581--84}). He also rejected the plurality's redressability argument, presuming that federal agencies would follow an authoritative Supreme Court interpretation of the ESA (\href{https://supreme.justia.com/cases/federal/us/504/555/}{504 U.S. at 585}). Nonetheless, Stevens concurred in reversal on the merits, applying the presumption against extraterritoriality to conclude that Congress did not intend \S~7(a)(2) to govern actions in foreign countries (\href{https://supreme.justia.com/cases/federal/us/504/555/}{504 U.S. at 585--89}).

\subsection{Dissenting Opinion (Justice Blackmun, joined by Justice O'Connor)}
Justice Blackmun argued that the majority misapplied the summary-judgment standard under Rule 56(c) by weighing evidence rather than asking whether a reasonable factfinder could conclude from Kelly's and Skilbred's past trips and wildlife-preservation careers that they would soon return (\href{https://supreme.justia.com/cases/federal/us/504/555/}{504 U.S. at 590--93}). He criticized the majority's demand for plane tickets or fixed return dates as ``an empty formality'' that resurrected ``code-pleading formalism'' (\href{https://supreme.justia.com/cases/federal/us/504/555/}{504 U.S. at 592--93}). On redressability, Blackmun accused the Government of playing ``Three-Card Monte'' by disavowing its prior view that the Secretary's regulations bound action agencies, noted that AID and the State Department had participated actively in the litigation (implicating collateral estoppel), and emphasized that AID's \$170 million contribution was substantial for Sri Lanka (\href{https://supreme.justia.com/cases/federal/us/504/555/}{504 U.S. at 595--601}). Finally, Blackmun warned that Part IV's rejection of procedural-injury standing reflected ``an anachronistically formal view of the separation of powers'' that improperly transferred power from Congress to the Executive and amounted to ``a slash-and-burn expedition through the law of environmental standing'' (\href{https://supreme.justia.com/cases/federal/us/504/555/}{504 U.S. at 602--06}).

\section{Examples: Future Applications}
\begin{enumerate}
\item \textbf{Same-Side (No Standing):} A wildlife photographer in Ohio who visited a rainforest in Brazil five years ago sues a federal export-credit agency for failing to conduct an environmental review of a Brazilian highway loan, stating in an affidavit that she ``hopes to return to Brazil someday once her schedule clears.'' \textit{Why:} Just as in \textit{Lujan}, past visits and indefinite ``some day'' intentions without concrete travel plans fail to establish an actual or imminent injury in fact (\href{https://supreme.justia.com/cases/federal/us/504/555/}{504 U.S. at 564}).
\item \textbf{Same-Side (No Standing):} Congress enacts a statute authorizing ``any U.S. citizen'' to sue the Department of Energy whenever the agency misses a statutory deadline to publish a technical compliance report, and a citizen who has no personal use for the report sues solely to compel agency obedience to the law. \textit{Why:} Under Part IV of \textit{Lujan}, Congress cannot manufacture Article III standing by conferring an abstract, self-contained procedural right to vindicate a generalized interest in executive compliance with federal law (\href{https://supreme.justia.com/cases/federal/us/504/555/}{504 U.S. at 573--77}).
\item \textbf{Opposite-Side (Standing):} A homeowner living directly downstream from a federally licensed hydroelectric dam sues the licensing agency for failing to prepare a statutorily required environmental impact statement before approving the dam. \textit{Why:} As Justice Scalia expressly noted in Footnote 7, a plaintiff enforcing a procedural right designed to protect a concrete, geographically linked interest has standing even without proving with certainty that the environmental study would halt the project (\href{https://supreme.justia.com/cases/federal/us/504/555/}{504 U.S. at 572 \& n.7}).
\item \textbf{Opposite-Side (Standing):} A marine biologist who holds a funded university grant, airline reservations, and research permits for a scheduled trip next month to study sea turtles at a specific coastal reef sues to enjoin a federal agency from dredging that exact reef without required consultation. \textit{Why:} Unlike Kelly and Skilbred, the researcher demonstrates a concrete, particularized, and imminent injury tied directly to the specific site affected by the challenged federal action (\href{https://supreme.justia.com/cases/federal/us/504/555/}{504 U.S. at 564, 566--67}).
\item \textbf{Fence-Sitter:} A federal statute grants every consumer a statutory right to receive accurate credit-report disclosures and authorizes statutory damages for violations; a consumer whose credit file contains an inaccurate internal alert that was never disseminated to any third party sues the credit bureau. \textit{Why:} Under Justice Kennedy's concurrence, Congress may ``define injuries and articulate chains of causation'' where none existed at common law (\href{https://supreme.justia.com/cases/federal/us/504/555/}{504 U.S. at 580}), suggesting standing might exist; yet under Scalia's strict concrete-injury requirement in Part IV, a bare statutory violation without real-world harm or a close common-law analog risks failing Article III (a tension that later divided the Court in \textit{Spokeo, Inc. v. Robins}, 578 U.S. 330 (2016), and \textit{TransUnion LLC v. Ramirez}, 594 U.S. 413 (2021)).
\end{enumerate}

\section{Critique}
\textbf{Originalist and Textualist Perspectives:} Defenders of \textit{Lujan}---building on Justice Scalia's influential 1983 article, \textit{The Doctrine of Standing as an Essential Element of the Separation of Powers}, 17 Suffolk U. L. Rev. 881---praise the decision for cabining unelected federal judges to traditional private-rights adjudication and preventing Congress from deputizing private attorneys general to supervise executive law enforcement (\href{https://www.fjc.gov/history/cases/cases-that-shaped-the-federal-courts/lujan-v-defenders-wildlife}{Federal Judicial Center}; \href{https://fedsoc.org/fedsoc-review/transunion-article-iii-and-expanding-the-judicial-role}{Federalist Society Review}). Yet \textit{Lujan} has also drawn sharp criticism from originalist and historical scholars. Scholars examining eighteenth-century English and early American practice note that the ``injury-in-fact'' formula was a twentieth-century doctrinal invention rather than a Founding-era requirement, and that early Congresses enacted informer and \textit{qui tam} statutes authorizing private citizens to sue without traditional personal injury (\href{https://lawreview.gmu.edu/print__issues/charting-a-course-past-spokeo-and-transunion/}{George Mason Law Review}; \href{https://constitution.congress.gov/browse/essay/artIII-S2-C1-6-2/ALDE_00012993/}{Constitution Annotated}). Justice Thomas and several originalist scholars have therefore urged grounding Article III limits in the historical distinction between private rights (where any legal invasion traditionally supported suit) and public rights (which generally required special damage) rather than \textit{Lujan}'s judge-made ``injury-in-fact'' test (\href{https://lawreview.gmu.edu/print__issues/charting-a-course-past-spokeo-and-transunion/}{George Mason Law Review}).

\textbf{Progressive and Doctrinal Perspectives:} In his landmark critique, Cass R. Sunstein argued that \textit{Lujan}'s invalidation of a congressionally authorized citizen suit lacked support in the text or history of Article III, privileging common-law private interests over modern regulatory beneficiaries (\href{https://chicagounbound.uchicago.edu/context/journal_articles/article/12523/viewcontent/1458.pdf}{Sunstein, 91 Mich. L. Rev. 163 (1992)}; \href{https://scholarship.law.ua.edu/cgi/viewcontent.cgi?article=1107&context=fac_articles}{Elliott, Stan. L. Rev.}). Similarly, Richard J. Pierce, Jr., contended that \textit{Lujan} imposed an unjustified judicial limit on Congress's power to create legally enforceable statutory rights and criticized the Court's ``plane ticket'' imminence rule as formalistic and manipulable (\href{https://scholarship.law.duke.edu/cgi/viewcontent.cgi?article=3222&context=dlj}{Pierce, 42 Duke L.J. 1170 (1993)}). Environmental scholars further argue that \textit{Lujan} systematically disadvantages diffuse ecological harms and ``existence-value'' interests---where the loss of a species harms millions of people subtly rather than one property owner acutely (\href{https://harvardlawreview.org/print/vol-129/existence-value-standing/}{Harvard Law Review}).

\textbf{Underlying Competing Values:} \textit{Lujan} reflects a fundamental clash between two visions of separation of powers and administrative accountability: (1) a formalist, executive-centered model that restricts federal courts to individualized private disputes to protect executive discretion and majoritarian politics from judicial overreach, versus (2) a legislative-supremacy and rule-of-law model that empowers Congress to enlist citizen suits to ensure that executive agencies faithfully obey statutory mandates.

\section{Key Quotations}
\begin{quote}
``Over the years, our cases have established that the irreducible constitutional minimum of standing contains three elements. First, the plaintiff must have suffered an `injury in fact'---an invasion of a legally protected interest which is (a) concrete and particularized, ... and (b) `actual or imminent, not ``conjectural'' or ``hypothetical,'' ' ... Second, there must be a causal connection between the injury and the conduct complained of---the injury has to be `fairly ... trace[able] to the challenged action of the defendant, and not ... th[e] result [of] the independent action of some third party not before the court.' ... Third, it must be `likely,' as opposed to merely `speculative,' that the injury will be `redressed by a favorable decision.' ''
\end{quote}
\noindent --- Justice Scalia, Opinion of the Court, 504 U.S. at 560--61 (\href{https://supreme.justia.com/cases/federal/us/504/555/}{Justia}; \href{https://www.courtlistener.com/opinion/112747/lujan-v-defenders-of-wildlife/}{CourtListener}).

\begin{quote}
``Such `some day' intentions---without any description of concrete plans, or indeed even any specification of when the some day will be---do not support a finding of the `actual or imminent' injury that our cases require.''
\end{quote}
\noindent --- Justice Scalia, Opinion of the Court, 504 U.S. at 564 (\href{https://supreme.justia.com/cases/federal/us/504/555/}{Justia}; \href{https://www.courtlistener.com/opinion/112747/lujan-v-defenders-of-wildlife/}{CourtListener}).

\begin{quote}
``To permit Congress to convert the undifferentiated public interest in executive officers' compliance with the law into an `individual right' vindicable in the courts is to permit Congress to transfer from the President to the courts the Chief Executive's most important constitutional duty, to `take Care that the Laws be faithfully executed,' Art. II, \S~3.''
\end{quote}
\noindent --- Justice Scalia, Opinion of the Court, 504 U.S. at 577 (\href{https://supreme.justia.com/cases/federal/us/504/555/}{Justia}; \href{https://www.courtlistener.com/opinion/112747/lujan-v-defenders-of-wildlife/}{CourtListener}).

\begin{quote}
``In my view, Congress has the power to define injuries and articulate chains of causation that will give rise to a case or controversy where none existed before, and I do not read the Court's opinion to suggest a contrary view.''
\end{quote}
\noindent --- Justice Kennedy, concurring in part and concurring in the judgment, 504 U.S. at 580 (\href{https://supreme.justia.com/cases/federal/us/504/555/}{Justia}; \href{https://www.courtlistener.com/opinion/112747/lujan-v-defenders-of-wildlife/}{CourtListener}).

\begin{quote}
``In conclusion, I cannot join the Court on what amounts to a slash-and-burn expedition through the law of environmental standing.''
\end{quote}
\noindent --- Justice Blackmun, dissenting, 504 U.S. at 606 (\href{https://supreme.justia.com/cases/federal/us/504/555/}{Justia}; \href{https://www.courtlistener.com/opinion/112747/lujan-v-defenders-of-wildlife/}{CourtListener}).


\end{document}
